% =====================================================================
% KHẢO LUẬN VỀ HỆ THUẬT NGỮ CỦA BẢN CHUYỂN NGỮ
% Phần do người chuyển ngữ biên soạn, không thuộc nguyên tác NCTM.
% =====================================================================

\cleardoublepage
\gdef\CurrentRunningHead{Khảo luận về hệ thuật ngữ của bản chuyển ngữ}

\refstepcounter{parallelstarchapter}
\bookmarksetup{startatroot}
\phantomsection
\addcontentsline{toc}{chapter}{Khảo luận về hệ thuật ngữ của bản chuyển ngữ}

\begingroup
\clubpenalty=10000
\widowpenalty=10000
\displaywidowpenalty=10000

\thispagestyle{plain}

\begin{center}
{\small\bfseries KHẢO LUẬN CỦA BẢN CHUYỂN NGỮ\par}

\vspace{0.9\baselineskip}

{\Large\bfseries Khảo luận về hệ thuật ngữ\par}

\vspace{0.15\baselineskip}

{\Large\bfseries của bản chuyển ngữ\par}

\vspace{0.65\baselineskip}

{\large Nguyên tắc phân biệt khái niệm và các quyết định chuyển ngữ trọng yếu\par}
\end{center}

\vspace{1.15\baselineskip}

\begingroup
\small

\noindent
\textit{Ghi chú của người chuyển ngữ.}
Phần khảo luận này do người chuyển ngữ biên soạn, không thuộc nguyên tác \emph{Focus in High School Mathematics: Reasoning and Sense Making}. Mục đích của nó là giải trình hệ thuật ngữ được sử dụng trong bản tiếng Việt: không chỉ nêu kết quả cuối cùng, mà còn trình bày những phân biệt khái niệm, những phương án từng được cân nhắc, cách các phương án ấy được kiểm định trên toàn văn, và lí do dẫn đến lựa chọn cuối cùng. Vì vậy, đây cần được đọc như một phần hậu khảo của bản dịch, chứ không phải như lời diễn giải thay cho quan niệm của NCTM.

\endgroup

\clearpage

\section*{1. Dịch một hệ khái niệm, không dịch những từ đứng riêng}

Khó khăn trung tâm của việc chuyển ngữ cuốn sách này không nằm ở chỗ tìm một từ tiếng Việt có nghĩa gần với một từ tiếng Anh. Khó khăn nằm ở chỗ nguyên tác sử dụng một \emph{hệ} khái niệm có nhiều nút nằm rất gần nhau về nghĩa nhưng không đồng nhất: \emph{reasoning}, \emph{inference}, \emph{argument}, \emph{argumentation}, \emph{justification}, \emph{proof}, \emph{explanation}; và ở một phía khác là \emph{sense making}, \emph{understanding}, \emph{meaning}, \emph{interpretation}. Nếu mỗi từ được dịch riêng theo thói quen từ điển, bản dịch có thể vẫn trôi chảy ở từng câu nhưng làm sụp đổ những đường biên khái niệm mà cuốn sách cố ý duy trì.

Bởi vậy, nguyên tắc được đặt ra trong giai đoạn hiệu chỉnh cuối không phải là ``một từ tiếng Anh tương ứng với một từ tiếng Việt'', mà là: \emph{mỗi lựa chọn phải được kiểm tra trong quan hệ với toàn bộ những từ lân cận của nó}. Một phương án chỉ được chấp nhận khi nó vừa đúng trong câu đang xét, vừa không chiếm mất chỗ cần dành cho một khái niệm khác ở những chương khác.

Cách làm này dẫn đến một thay đổi quan trọng so với những bản dịch thử nghiệm ban đầu. Có những phương án từng tỏ ra rất hấp dẫn khi nhìn riêng lẻ nhưng không đứng vững khi đặt vào toàn bộ văn bản. Chẳng hạn, \emph{reasoning} từng được thử bằng ``lập luận'', ``luận'' và ``suy luận''. Cả ba đều có căn cứ ngôn ngữ nhất định. Tuy nhiên, khi đọc đồng thời những đoạn mà NCTM dùng \emph{reasoning} bên cạnh \emph{argument}, hoặc dùng \emph{statistical reasoning} bên cạnh \emph{statistical inference}, ta buộc phải dành ``lập luận'' cho \emph{argument} và ``suy luận'' cho \emph{inference}. Chính việc đọc theo mạng lưới như vậy mở đường cho lựa chọn cuối cùng \emph{reasoning} là ``suy lí''.

Quá trình kiểm định được thực hiện theo hai tầng. Tầng thứ nhất là đọc và lập bản đồ các họ từ trong toàn văn: dạng danh từ, động từ, các tổ hợp chuyên môn, những nơi hai hay nhiều thuật ngữ cùng xuất hiện. Tầng thứ hai là kiểm tra ``xung đột'': nếu hai từ tiếng Anh khác nhau có nguy cơ bị dồn vào cùng một từ tiếng Việt, phải đọc lại toàn bộ những ngữ cảnh có liên quan trước khi chốt. Nhờ đó, các cặp đặc biệt nhạy cảm như \emph{reasoning--inference}, \emph{reasoning--argument}, \emph{argument--argumentation}, \emph{justification--proof--explanation}, \emph{sense making--understanding--meaning}, và \emph{interpretation--explanation} được giữ tách biệt.

Điều cần nhấn mạnh là đây là một hệ thuật ngữ dành cho \emph{bản dịch này}. Nó không nhằm tuyên bố rằng tiếng Việt nói chung, trong mọi ngành và mọi thời kì, bắt buộc phải phân chia các từ theo đúng cách dưới đây. Nhiều từ trong số này có miền nghĩa giao nhau trong sử dụng tự nhiên. Việc chuyên biệt hóa ở đây là một quyết định dịch thuật có chủ đích, nhằm bảo toàn kiến trúc khái niệm của một văn bản giáo dục toán học cụ thể.

\section*{2. \emph{Reasoning}: vì sao chọn ``suy lí''}

NCTM đưa ra một định nghĩa khái quát: \emph{reasoning} có thể được xem là quá trình rút ra kết luận trên cơ sở bằng chứng hoặc những giả định đã nêu. Nhưng ngay sau đó, cuốn sách mở rộng phạm vi của khái niệm: suy lí toán học có thể đi từ giải thích và biện minh phi hình thức, quan sát quy nạp, tìm tòi, phỏng đoán, khởi đầu sai hướng, cho tới suy diễn hình thức và chứng minh. Vì thế, \emph{reasoning} trong cuốn sách không phải chỉ là một phép suy ra đơn lẻ; nó là hoạt động nhận thức rộng hơn, trong đó người học lần theo quan hệ, căn cứ, cấu trúc và hệ quả để đi từ cái đã biết tới điều có thể kết luận.

Phương án ``lập luận'' từng có sức hấp dẫn vì trong tiếng Việt nó gợi rõ hoạt động có lí lẽ. Tuy nhiên, nguyên tác đồng thời dùng \emph{argument} và \emph{argumentation}. Ở nhiều đoạn, học sinh \emph{reason}, rồi tinh chỉnh \emph{arguments} để có thể truyền đạt chúng; ở chỗ khác, \emph{formal argumentation} là hoạt động xây dựng và thẩm định các chuỗi lí lẽ. Nếu \emph{reasoning} cũng được gọi là ``lập luận'', sự khác nhau giữa hoạt động tư duy rộng và sản phẩm lập luận có cấu trúc sẽ bị xóa.

Phương án ``suy luận'' cũng từng được xem là ứng viên mạnh. Vấn đề chỉ lộ rõ khi sang thống kê. NCTM phân biệt \emph{statistical reasoning} với \emph{statistical inference}; hơn nữa, trong cùng một câu có thể xuất hiện cả \emph{inferential reasoning} và \emph{scope of inference}. Nếu \emph{reasoning} được dịch là ``suy luận'', ta không còn một từ rõ ràng để giữ riêng \emph{inference}. Đây không phải khác biệt hình thức. Suy lí thống kê bao gồm cả việc đặt câu hỏi, xem xét biến thiên, lựa chọn biểu diễn, đánh giá tính hợp lí, suy nghĩ dưới điều kiện bất định; còn suy luận thống kê chỉ một lớp thao tác đặc thù trong đó từ dữ liệu mẫu người ta đi tới kết luận rộng hơn, với phạm vi suy luận bị quy định bởi cách dữ liệu được thu thập.

Phương án ``luận'' có ưu điểm cô đọng và từng được khai thác như một từ Hán Việt giàu khả năng tạo nghĩa. Nhưng khi đi vào câu văn toàn sách, các cấu trúc như ``học sinh luận'', ``thói quen luận'', ``luận với hàm số'' dễ trở nên quá nén, đồng thời vẫn đứng rất gần ``lập luận''. Sau khi cân nhắc toàn hệ, bản dịch chọn ``suy lí''.

Trong lựa chọn này, ``lí'' được hiểu theo trường nghĩa của căn cứ, quan hệ, cấu trúc và điều khiến một kết luận đứng vững. Có thể neo khái niệm bằng một câu: \emph{suy lí là hoạt động của tư duy trong đó ta lần theo những cái lí --- quan hệ, cấu trúc, quy luật, căn cứ --- của những gì đã biết để đi tới nhận định hoặc kết luận mới}. Câu này là một cách giải nghĩa phục vụ bản dịch, không phải một định nghĩa từ nguyên có tham vọng áp dụng cho mọi cách dùng tiếng Việt.

Việc chọn chữ ``lí'' còn có một điểm tựa đáng chú ý trong lịch sử nghĩa của chữ {\hanfont 理}. \emph{Thuyết văn giải tự} giảng {\hanfont 理} là ''trị ngọc'' ({\hanfont 治玉}), tức chế tác, mài giũa ngọc; thành tố {\hanfont 玉} trong chữ cũng cho thấy mối liên hệ ban đầu ấy. Từ việc nhận ra và thuận theo những đường phân chia, đường vân và kết cấu của vật chất khi chế tác, {\hanfont 理} còn được dùng để chỉ thớ, vân hay kết cấu của sự vật, như trong {\hanfont 木理} --- thớ hoặc vân gỗ, {\hanfont 肌理} --- thớ hay kết cấu của cơ thể, và {\hanfont 紋理} --- đường vân, kết cấu. Từ trường nghĩa cụ thể này, {\hanfont 理} mở rộng sang những nghĩa trừu tượng hơn: sự phân định, trật tự, quy luật, nguyên tắc, cái khiến sự vật có cấu trúc và có thể được nhận biết theo một mạch nhất định. Sự chuyển nghĩa ấy không phải là căn cứ duy nhất để dịch \emph{reasoning} thành ``suy lí'', nhưng nó giúp làm hiện rõ sức gợi của chữ ``lí'': suy lí có thể được hình dung như việc lần theo những đường nét và quan hệ vốn làm nên cấu trúc của sự vật, thay vì chỉ thực hiện một phép suy ra hình thức.

Từ lựa chọn hạt nhân ấy, một họ thuật ngữ trở nên nhất quán: \emph{mathematical reasoning} là ``suy lí toán học''; \emph{formal reasoning} là ``suy lí hình thức''; \emph{statistical reasoning} là ``suy lí thống kê''; \emph{geometric reasoning} là ``suy lí hình học''; \emph{empirical reasoning} là ``suy lí thực nghiệm''; \emph{deductive reasoning} là ``suy lí suy diễn''; \emph{inductive reasoning} là ``suy lí quy nạp''; \emph{reasoning habits} là ``thói quen suy lí''.

Trong hệ van Hiele, \emph{relational-inferential reasoning} được diễn đạt thành ``suy lí dựa trên quan hệ và suy luận'', tránh một tổ hợp máy móc vừa khó đọc vừa làm mờ hai lớp khái niệm.

Không phải mọi dạng của từ \emph{reason} đều buộc phải mang ``suy lí''. \emph{Reason} khi là danh từ thường là ``lí do'' hoặc ``căn cứ''. \emph{Reasonable} và \emph{reasonableness} trong các ngữ cảnh đánh giá kết quả được chuyển thành ``hợp lí'' và ``tính hợp lí''. Dạng \emph{reasoned} cũng được xử lí theo cú pháp: \emph{reasoned solving} là ``giải dựa trên suy lí'', còn \emph{reasoned connections} có thể là những kết nối ``được thiết lập bằng suy lí''. Như vậy, bảng thuật ngữ chỉ định hướng cách hiểu; việc lựa chọn cách dịch cụ thể vẫn phải dựa vào ngữ cảnh của từng câu.

\section*{3. \emph{Sense making}: từ ``hiểu'' đến ``kiến tạo ý nghĩa''}

Nếu chỉ đọc tiêu đề \emph{Reasoning and Sense Making}, dịch \emph{sense making} là ``hiểu'' có vẻ tự nhiên. Chính bản dịch ở một giai đoạn trước đã đi theo hướng đó. Tuy nhiên, định nghĩa của NCTM tạo ra một phép thử quyết định: \emph{sense making} được định nghĩa là việc \emph{developing understanding} về một tình huống, bối cảnh hoặc khái niệm bằng cách nối nó với tri thức đã có. Nói cách khác, \emph{sense making} là hoạt động đang diễn ra; \emph{understanding} là cái được phát triển nhờ hoạt động ấy.

Nếu cả hai cùng được dịch bằng ``hiểu'', cấu trúc định nghĩa bị làm phẳng: ``hiểu là phát triển sự hiểu''. Vấn đề không chỉ là câu văn vụng, mà là mất phân biệt giữa quá trình và kết quả. Bản dịch cuối cùng vì thế chọn:

\begin{quote}
\emph{Chúng tôi định nghĩa kiến tạo ý nghĩa là việc phát triển sự hiểu biết về một tình huống, bối cảnh hoặc khái niệm bằng cách kết nối nó với tri thức đã có.}
\end{quote}

``Kiến tạo ý nghĩa'' cố giữ ba đặc điểm của khái niệm. Thứ nhất, đây là một hoạt động chứ không chỉ là trạng thái. Thứ hai, hoạt động ấy hướng tới việc làm cho đối tượng trở nên có nghĩa đối với người học. Thứ ba, ý nghĩa không được tạo ra trong khoảng trống mà hình thành qua những kết nối với tri thức và kinh nghiệm đã có. Từ ``kiến tạo'' ở đây không nhằm gán toàn bộ một học thuyết kiến tạo luận cho NCTM; nó chỉ làm nổi bật tính chủ động và tính quan hệ đã có ngay trong định nghĩa của nguyên tác.

Một số phương án khác cũng đã được cân nhắc. ``Quá trình hiểu'' hay ``tiến trình hiểu'' vẫn giữ từ ``hiểu'' ở trung tâm và do đó chưa giải quyết được xung đột với \emph{understanding}. ``Lập nghĩa'' quá dễ gợi một hành vi đặt hoặc xác lập nghĩa theo cách chủ ý, và trong một số truyền thống ngôn ngữ học hoặc lôgic còn có thể gợi những liên tưởng không thuộc ngữ cảnh này. ``Diễn nghĩa'' lại gần với việc diễn giải một ý nghĩa đã có hơn là quá trình hình thành sự hiểu biết bằng kết nối. ``Kiến tạo ý nghĩa'' vì thế là phương án giữ được nhiều nhất cấu trúc khái niệm của NCTM.

Tuy nhiên, lựa chọn thuật ngữ không có nghĩa mọi động từ \emph{make sense of} đều phải được thay cơ học bằng ``kiến tạo ý nghĩa cho''. Trong câu văn tự nhiên, tùy chức năng, có thể dùng ``hiểu'', ``hiểu được'', ``nhận ra ... có nghĩa gì'', hoặc ``làm cho ... trở nên có nghĩa''. Chỉ khi cụm từ thực sự mang vai trò của khái niệm kĩ thuật đang được bàn tới thì ``kiến tạo ý nghĩa'' cần được giữ rõ. Đây là một ví dụ quan trọng cho nguyên tắc: nhất quán khái niệm không đồng nghĩa với nhất loạt bề mặt.

\section*{4. \emph{Understanding}, \emph{meaning} và \emph{interpretation}}

Sau khi tách \emph{sense making} khỏi \emph{understanding}, trường nghĩa ``hiểu'' trở nên sáng hơn. \emph{Understanding} khi là danh từ được ưu tiên dịch là ``sự hiểu biết''; \emph{understand} là ``hiểu''. Vì thế \emph{conceptual understanding} trong mô hình năng lực toán học của \emph{Adding It Up} là ``sự hiểu biết khái niệm''. Cách dịch này cho phép NCTM nói một cách tự nhiên rằng kiến tạo ý nghĩa \emph{phát triển sự hiểu biết}, và rằng một chứng minh tốt có thể giúp người học \emph{hiểu} ý nghĩa của điều được chứng minh.

\emph{Meaning} được giữ là ``ý nghĩa''. Các tính từ \emph{meaningful} và \emph{meaningless} được dịch tùy cấu trúc thành ``có nghĩa'', ``có ý nghĩa'', ``vô nghĩa'' hoặc ``không có ý nghĩa''. Sự phân biệt này cho phép ta nói: kiến tạo ý nghĩa là việc làm; ý nghĩa là cái dần hiện ra trong mạng lưới quan hệ; sự hiểu biết là cái hình thành và phát triển ở người học nhờ những quan hệ ấy.

\emph{Interpret} và \emph{interpretation} được dành cho ``diễn giải''. Trong cuốn sách, đặc biệt ở thống kê và giải quyết vấn đề, diễn giải là xác định một kết quả, một mô hình hay một con số \emph{có nghĩa gì trong bối cảnh đang xét} và nó trả lời câu hỏi ban đầu ra sao. Vì vậy, ``diễn giải'' phải được phân biệt với ``giải thích''. Một lời giải thích thường trả lời vì sao hoặc bằng cách nào một điều xảy ra hay đúng; một sự diễn giải làm rõ ý nghĩa của kết quả trong ngữ cảnh. Cũng vì lí do ấy, \emph{mathematical exposition} được chuyển thành ``trình bày toán học'', không dùng ``diễn giải toán học'', để không chiếm mất từ dành cho \emph{interpretation}.

\emph{Comprehension} xuất hiện thưa hơn và không tạo thành một nút khái niệm độc lập đủ mạnh trong cuốn sách. Vì vậy bản dịch không cưỡng ép một thuật ngữ riêng chỉ để tạo đối xứng từ điển; nó được xử lí theo câu bằng ``sự hiểu'', ``khả năng hiểu'' hoặc cách diễn đạt phù hợp với đối tượng được nói tới.

\section*{5. \emph{Inference}: dành ``suy luận'' cho thao tác suy ra}

Việc chọn ``suy lí'' cho \emph{reasoning} cho phép giữ ``suy luận'' cho \emph{inference}. Sự phân biệt này đặc biệt quan trọng trong Chương 8. \emph{Statistical reasoning} là ``suy lí thống kê'', còn \emph{statistical inference} là ``suy luận thống kê''.

Hai cụm có quan hệ gần nhưng không đồng nhất. Suy lí thống kê là trường hoạt động rộng: đọc dữ liệu, đặt câu hỏi, nhận diện biến thiên, chọn cách biểu diễn, xem xét xác suất, đánh giá tính hợp lí, suy nghĩ trong điều kiện bất định. Suy luận thống kê là một lớp hoạt động đặc thù trong trường ấy, liên quan đến việc từ dữ liệu đi tới kết luận vượt ra ngoài chính dữ liệu quan sát, với phạm vi kết luận phụ thuộc vào thiết kế thu thập dữ liệu. Bởi vậy \emph{scope of inference} được dịch là ``phạm vi suy luận''.

Đối với \emph{inferential reasoning}, bản dịch tránh tổ hợp ``suy lí suy luận''. Trong ngữ cảnh thống kê, cách diễn đạt được tái cấu trúc thành ``suy lí trong suy luận thống kê''; \emph{informal inferential reasoning} là ``suy lí phi hình thức trong suy luận thống kê''. Cấu trúc này dài hơn tiếng Anh nhưng bảo toàn được cả hai tầng: \emph{reasoning} là hoạt động nhận thức, \emph{inference} là loại thao tác thống kê mà hoạt động ấy đang phục vụ.

Tương tự, \emph{deduction} và \emph{induction} không bị nhập vào \emph{inference}. \emph{Deduction} là ``suy diễn'', còn khi nằm trong tổ hợp \emph{deductive reasoning} thì thành ``suy lí suy diễn''. \emph{Inductive reasoning} là ``suy lí quy nạp''. Việc dành những từ riêng này giúp các chương về tiến trình phát triển của suy lí và hình học giữ được sự phân tầng từ thực nghiệm, quy nạp, phi hình thức tới hình thức.

\section*{6. \emph{Argument} và \emph{argumentation}: sản phẩm và hoạt động}

Trong ngôn ngữ thường ngày, ``lập luận'' đôi khi được dùng rộng tới mức gần với ``suy lí''. Bản dịch này cố ý thu hẹp nó để tương ứng với \emph{argument}. Một \emph{argument} là một cấu trúc lí lẽ có thể được trình bày, đọc, đánh giá, tinh chỉnh hoặc phản biện. Nó có tính chất của một sản phẩm được tổ chức và có thể giao tiếp. Trong khi đó, \emph{reasoning} rộng hơn: nó có thể diễn ra trước khi một lập luận hoàn chỉnh được hình thành.

Sự phân biệt ấy hiện rõ khi NCTM nói học sinh phải có khả năng xác định tính hợp lệ của một \emph{argument}, hoặc phải \emph{refine arguments so that they can be effectively communicated}. Ở những chỗ ấy, ``lập luận'' chính xác hơn ``suy lí'', vì đối tượng đang được thẩm định hoặc chỉnh sửa là một cấu trúc lí lẽ đã có hình thức tương đối rõ.

\emph{Argumentation} lại chỉ hoạt động tạo lập, trao đổi và thẩm định lập luận. Vì thế nó được chuyển là ``hoạt động lập luận'', và \emph{formal argumentation} là ``hoạt động lập luận hình thức''. Sự thêm vào của từ ``hoạt động'' không phải phần tô điểm; nó giữ phân biệt ngữ pháp--khái niệm giữa \emph{argument} và \emph{argumentation}.

\emph{Partial arguments} được dịch là ``lập luận từng phần''. Phương án ``lập luận bộ phận'' từng có thể nghĩ tới, nhưng ``từng phần'' diễn đạt tự nhiên hơn tình trạng một lập luận mới chỉ hoàn thành một phần, còn thiếu những mắt xích cần thiết để trở thành lập luận đầy đủ.

Hình 10.1 cung cấp một phép thử nhỏ nhưng rõ: rubric đánh giá câu trả lời dùng \emph{argument}, vì thế phải là ``lập luận'', không phải ``biện giải''. \emph{A clear argument with minimal gaps} được hiểu là ``một lập luận rõ ràng với rất ít chỗ thiếu sót''. Sự nhất quán ở những chi tiết như vậy giúp hệ thuật ngữ không chỉ đúng trong chương lí thuyết mà còn sống được trong những nhiệm vụ đánh giá cụ thể.

\section*{7. \emph{Justification}, \emph{proof} và \emph{explanation}}

Ba khái niệm này có thể cùng trả lời câu hỏi ``vì sao'', nhưng chúng không
đồng nhất.

\emph{Justify} và \emph{justification} được dịch là ``biện minh'' và ``sự biện minh'' khi cần danh hóa. Trong văn nói, ``biện minh'' đôi khi có sắc thái phòng vệ; trong ngữ cảnh nhận thức của cuốn sách, nó được dùng theo nghĩa đưa ra căn cứ đủ sức làm cho một mệnh đề, một lựa chọn hay một lời giải được chấp nhận. Chính Chương 2 cung cấp một ranh giới rất có giá trị: bằng chứng thực nghiệm có thể \emph{support} một phỏng đoán nhưng không \emph{justify} nó. Vì vậy ``hỗ trợ'' và ``biện minh'' không thể đổi chỗ cho nhau.

\emph{Proof} được giữ là ``chứng minh''; khi cần nhấn mạnh sản phẩm, có thể nói ``một chứng minh''. \emph{Formal proof} là ``chứng minh hình thức''. \emph{Proving}, khi câu văn thực sự nhấn vào quá trình, được diễn đạt là ``hoạt động chứng minh'' hoặc ``quá trình chứng minh''. Bản dịch tránh biến mọi \emph{proof} thành một cách nói về \emph{reasoning}, bởi chứng minh là một kết quả hay hình thức trình bày có yêu cầu riêng. Câu mô tả vai trò của chứng minh được dịch theo hướng: ``Chứng minh là một cách trình bày suy lí hình thức được xây dựng trên nền tảng kiến tạo ý nghĩa''. Cách nói này giữ quan hệ giữa hai khái niệm mà không đồng nhất chúng.

\emph{Explanation} là ``giải thích''. \emph{Informal explanation} là ``giải thích phi hình thức''; \emph{partial explanations} tùy câu được chuyển thành ``giải thích từng phần'' hoặc ``giải thích còn dang dở''. Một chứng minh có thể có chức năng giải thích --- nghĩa là giúp ta thấy không chỉ rằng một mệnh đề đúng mà còn vì sao nó đúng --- nhưng từ đó không suy ra \emph{proof} và \emph{explanation} là cùng một khái niệm. Chức năng có thể giao nhau trong khi loại hoạt động và chuẩn mực chấp nhận vẫn khác.

Do đó, trong bản dịch này có thể hình dung một chuỗi nhưng không phải một thang đồng nhất: người học suy lí; từ suy lí có thể hình thành những giải thích và lập luận; các lập luận có thể chứa sự biện minh; trong những bối cảnh đòi hỏi chắc chắn toán học, hoạt động lập luận hình thức có thể đi tới chứng minh. Những đường đi ấy giao nhau nhưng không cho phép thay tên tùy tiện.

\section*{8. Các tổ hợp \emph{reasoning} trong cấu trúc cuốn sách}

Việc chọn ``suy lí'' phải đứng vững không chỉ ở Chương 1 mà trên toàn cấu trúc cuốn sách. Năm chương nội dung vì thế mang các tên: ``Suy lí với số và đo lường'', ``Suy lí với kí hiệu đại số'', ``Suy lí với hàm số'', ``Suy lí với hình học'', và ``Suy lí với thống kê và xác suất''.

Cách dịch \emph{Reasoning with Algebraic Symbols} thành ``Suy lí với kí hiệu đại số'' là có chủ ý. Nguyên tác không nói chung về \emph{reasoning with algebra}; nó tập trung vào vai trò của biến, biểu thức, phương trình, dấu bằng và các phép biến đổi như những phương tiện biểu thị và thao tác trên quan hệ. Giữ ``kí hiệu đại số'' vì thế quan trọng hơn một tiêu đề Việt hóa ngắn gọn nhưng rộng nghĩa hơn.

Trong Chương 5, \emph{Reasoned Solving} được dịch là ``Giải dựa trên suy lí''. Phương án cũ ``Giải có cơ sở luận'' vừa giữ dấu vết của thuật ngữ đã bỏ, vừa khó đọc như một tên yếu tố then chốt. ``Dựa trên suy lí'' làm rõ điều NCTM muốn đối lập với việc thao tác phương trình như chuỗi thủ tục không có căn cứ. \emph{Mindful Manipulation} được giữ là ``Biến đổi có ý thức'', nhấn mạnh việc thao tác kí hiệu phải được dẫn dắt bởi mục tiêu và sự hiểu biết chứ không phải phản xạ thủ tục.

Ở Chương 7, các mức van Hiele được đọc như các mức \emph{suy lí}, không phải những loại ``lập luận'' đã hoàn chỉnh. Vì thế có ``suy lí trực quan--toàn thể'', ``suy lí mô tả--phân tích'', ``suy lí thực nghiệm'', ``suy lí dựa trên quan hệ và suy luận'', và ``suy lí hình thức''. Cách gọi này nối trực tiếp chương hình học với tiến trình phát triển của suy lí đã trình bày ở Chương 2.

\section*{9. Năm mạch năng lực toán học trong \emph{Adding It Up}}

Hình 2.1 đưa vào một hệ thuật ngữ có nguồn từ \emph{Adding It Up}. Đây là một ví dụ cho thấy không nên dịch từng nhãn độc lập mà phải nhìn cả hệ năm \emph{strands of mathematical proficiency}.

\emph{Conceptual understanding} được dịch là ``sự hiểu biết khái niệm'', thay cho ``hiểu khái niệm'', để giữ \emph{understanding} như một danh từ chỉ trạng thái hay năng lực đã phát triển. \emph{Procedural fluency} là ``sự thành thạo thủ tục''; \emph{strategic competence} là ``năng lực chiến lược''; \emph{adaptive reasoning} là ``suy lí thích ứng''; và \emph{productive disposition} là ``khuynh hướng tích cực''.

Lựa chọn ``khuynh hướng'' thay cho ``thái độ'' là có chủ ý. Nguyên tác giải thích \emph{productive disposition} bằng một \emph{habitual inclination}: một xu hướng bền vững nhìn toán học là có nghĩa, hữu ích và đáng học, đi cùng niềm tin vào sự chuyên cần và năng lực của chính mình. ``Thái độ tích cực'' không sai ở mức phổ thông, nhưng làm yếu tính bền vững và định hướng hành vi của \emph{disposition}.

Tương tự, phần giải thích \emph{adaptive reasoning} gồm \emph{logical thought, reflection, explanation, and justification}; bản dịch dùng ``năng lực tư duy lôgic, phản tư, giải thích và biện minh''. Việc dùng ``biện minh'' ở đây đồng bộ với hệ \emph{justification} đã xác lập ở trên.

\section*{10. \emph{Recursion} và \emph{recurrence}: vì sao ở Chương 3 là ``truy hồi''}

Một quyết định khác chỉ có thể được hiểu đúng khi xét ngữ cảnh là \emph{recursion}. Trong tiếng Anh, họ từ này bao quát nhiều hiện tượng. Khi trọng tâm là một hàm, thủ tục hay thuật toán tự gọi lại chính nó hoặc được xác định bằng sự tự quy chiếu, tiếng Việt ``đệ quy'' là lựa chọn tự nhiên: \emph{recursive function}, \emph{recursive algorithm}, \emph{recursive call} lần lượt là ``hàm đệ quy'', ``thuật toán đệ quy'', ``lời gọi đệ quy''.

Nhưng trong toán học phổ thông, nhất là khi nói về dãy, mẫu hình và quá trình rời rạc, \emph{recursive} thường mô tả quan hệ trong đó giá trị bước sau được xác định từ một hoặc nhiều giá trị bước trước. Ở trường hợp ấy, ``truy hồi'' làm nổi đúng cấu trúc bước sau--bước trước. Do đó có ``dãy truy hồi'', ``công thức truy hồi'', ``quy tắc truy hồi'', và \emph{recurrence relation} là ``hệ thức truy hồi''.

Trong Chương 3, \emph{recursion} xuất hiện như một nội dung nằm trong yếu tố ``Nhiều biểu diễn của hàm số''. Trọng tâm không phải một thuật toán tự gọi mà là cách biểu diễn một hàm số, dãy hoặc quá trình rời rạc bằng quy tắc xác định bước sau từ bước trước. Vì vậy bản dịch dùng ``truy hồi''. Đây là một ví dụ điển hình cho nguyên tắc: họ từ tiếng Anh có thể cần hơn một cách dịch tiếng Việt khi cấu trúc toán học được nói tới thay đổi.

\section*{11. Một số thuật ngữ chương trình và thống kê cần được giữ theo ngữ cảnh}

\emph{Coherence} là từ trung tâm của Chương 10. Tiêu đề chương được dịch ``Tính nhất quán''. Ở đây, ``nhất quán'' không chỉ nói các phát biểu không mâu thuẫn với nhau; nó chỉ sự liên thông của chương trình qua các khóa học, các mảng nội dung, giảng dạy và đánh giá, sao cho học sinh có thể nhìn toán học như một chỉnh thể. Khi nguyên tác dùng đồng thời \emph{coherent} và \emph{cohesive}, bản dịch có thể dùng ``nhất quán và gắn kết'' để giữ cả hai sắc thái. Phương án ``mạch lạc'' từng xuất hiện nhưng dễ gợi phẩm chất của văn bản hoặc lời nói hơn là cấu trúc liên thông của một chương trình học.

Trong Chương 9, \emph{tracking} và \emph{track} được giữ là ``phân luồng'' và ``luồng''. Đây là ngữ cảnh công bằng giáo dục, nơi học sinh được đưa vào những quỹ đạo hay cấp độ khóa học khác nhau. ``Phân ban'' chỉ một cấu trúc tổ chức khác và không bao quát đúng hiện tượng mà chương đang phân tích.

Trong thống kê, \emph{variation} được dịch là ``sự biến thiên''. Từ này không chỉ có nghĩa một đại lượng ``thay đổi'', mà chỉ hiện tượng các quan sát, mẫu hoặc kết quả khác nhau --- một đối tượng mà suy lí thống kê phải mô tả, mô hình hóa và đưa vào đánh giá. \emph{Association} trong ngữ cảnh dữ liệu được ưu tiên là ``mối liên hệ''. \emph{Sampling distribution} là ``phân phối lấy mẫu''; \emph{random selection} là ``chọn ngẫu nhiên''; \emph{random assignment} là ``phân bổ ngẫu nhiên''; và \emph{scope of inference}, như đã nói, là ``phạm vi suy luận''.

Cụm \emph{interpreting designed statistical studies} được hiểu theo hướng ``diễn giải các nghiên cứu thống kê có thiết kế''. ``Diễn giải'' ở đây là đọc kết quả trong quan hệ với cách nghiên cứu được thiết kế và với phạm vi kết luận mà thiết kế ấy cho phép. Điều đó tiếp tục cho thấy vì sao \emph{interpretation} phải được giữ khác \emph{explanation}.

\section*{12. Về chính tả và tính nhất quán hình thức}

Hệ thuật ngữ không thể tách khỏi chính tả. Bản dịch này nhất quán ưu tiên các dạng ``lí'', ``kĩ'', ``mĩ'', ``tỉ'' trong những trường hợp thuộc quy ước đã chọn; vì thế có ``suy lí'', ``hợp lí'', ``kĩ năng'', ``tỉ lệ''. Việc giữ ``i'' không phải một quyết định khái niệm như các phân biệt ở trên, nhưng có ý nghĩa biên tập: một hệ thuật ngữ chỉ thực sự ổn định khi hình thức chính tả cũng không dao động giữa các chương.

Tương tự, bản dịch tránh biến kí hiệu hoặc cấu trúc LaTeX thành một lớp diễn đạt thay cho văn xuôi nếu không cần thiết. Những quan hệ suy diễn, tương đương và kết nối thường được viết thành câu khi câu văn có thể biểu đạt rõ hơn. Mục tiêu chung là để thuật ngữ toán học chính xác mà văn bản vẫn có thể được đọc như tiếng Việt, không như một bảng thay thế từng từ từ tiếng Anh.

\section*{13. Kết luận: bảo toàn những đường biên cần thiết}

Kết quả quan trọng nhất của quá trình nghiên cứu thuật ngữ không phải là một danh sách song ngữ, mà là một \emph{bản đồ đường biên}. Trong bản đồ ấy, ``suy lí'' không phải ``suy luận''; suy lí cũng không phải ``lập luận''. ``Lập luận'' không phải ``hoạt động lập luận''. ``Biện minh'' không phải ``chứng minh'', và chứng minh không phải ``giải thích'', dù một chứng minh tốt có thể có sức giải thích. ``Kiến tạo ý nghĩa'' không phải ``sự hiểu biết''; nó là hoạt động có thể làm sự hiểu biết phát triển. ``Diễn giải'' không phải ``giải thích''; nó làm rõ kết quả có nghĩa gì trong bối cảnh. Những từ này được nối với nhau bởi quan hệ, chứ không bị hòa vào nhau.

Vì thế, các lựa chọn cốt lõi của bản chuyển ngữ có thể được đọc như một hệ: \emph{reasoning} là ``suy lí''; \emph{inference} là ``suy luận''; \emph{argument} là ``lập luận''; \emph{argumentation} là ``hoạt động lập luận''; \emph{justification} là ``biện minh''; \emph{proof} là ``chứng minh''; \emph{explanation} là ``giải thích''; \emph{sense making} là ``kiến tạo ý nghĩa''; \emph{understanding} là ``sự hiểu biết''; \emph{meaning} là ``ý nghĩa''; và \emph{interpretation} là ``diễn giải''. Mỗi từ chỉ thật sự được hiểu đầy đủ khi đặt vào quan hệ với những từ còn lại.

Sự lựa chọn này cũng giải thích vì sao bản dịch cuối cùng không giữ những cặp rất ngắn từng được thử nghiệm như ``luận và hiểu''. Sự cô đọng có giá trị, nhưng trong một văn bản mà chính đối tượng nghiên cứu là cách con người hiểu, suy nghĩ, rút ra kết luận, xây dựng lập luận và chứng minh, cái giá của việc làm mất một đường biên khái niệm là quá lớn. ``Suy lí và kiến tạo ý nghĩa'' dài hơn, nhưng nó cho phép toàn bộ mạng lưới còn lại có chỗ đứng rõ ràng.

Mục tiêu sau cùng của hệ thuật ngữ này không phải tạo ra một tiếng Việt xa lạ để đổi lấy vẻ chuyên môn. Mục tiêu là bảo toàn những phân biệt mà nguyên tác cần, đồng thời để bản dịch có thể sống như một văn bản giáo dục toán học bằng tiếng Việt. Khi hai yêu cầu ấy xung đột, bản dịch ưu tiên trung thành với cấu trúc khái niệm trước, rồi tìm cách làm câu văn tự nhiên trong phạm vi cấu trúc đó. Khi chúng không xung đột, tiếng Việt tự nhiên luôn được ưu tiên hơn lối dịch cơ học.

Đó cũng là nguyên tắc đã dẫn toàn bộ vòng hiệu chỉnh cuối: dịch NCTM cho đúng NCTM trước; chỉ sau khi hệ khái niệm của nguyên tác đã được xác lập rõ, những diễn giải, mở rộng hay kiến tạo thuật ngữ riêng mới có thể bắt đầu trên một nền tảng vững chắc.

\endgroup
